% !TEX root = ../enac_project_fr.tex
\chapter*{Introduction générale}
\label{introduction}
\addcontentsline{toc}{chapter}{\nameref{introduction}}
\markright{Introduction générale}
\thispagestyle{fancy}

Ce rapport présente un travail réalisé pour CS Group, qui est une entreprise française qui offre des services dans de nombreux domaines liés à la défense, l'aéronautique, le transport, l'énergie et le spatial.
Parmi ces différents domaine, l'entreprise est particulièrement active depuis plus de 40 ans dans le secteur du spatial, dont elle développe et fournit des solutions de surveillance spatiale, d'observation terrestre, de navigation et télécommunication.
Le stage s'inscrit dans le cadre des activités de recherche et développement en navigation par satellites, j'ai pour cela eu la chance d'être intégré à l'équipe FD NAV (\emph{Flight Dynamics and Navigation}) de la \emph{Buisiness Unit} Espace de l'entreprise.



\section{Présentation du contexte}

Depuis 2002, CS GROUP développe Orekit, une bibliothèque Java destinée à fournir des outils pour de nombreuses applications de dynamique du vol spatial, d'optimisation de trajectoires et d'autres aspects liés à la navigation.

Aujourd'hui, Orekit est largement utilisée dans le secteur spatial et est activement maintenue par les développeurs de CS GROUP ainsi que par une communauté d'utilisateurs. La bibliothèque est \textit{open source} et distribuée sous licence Apache 2.0, une licence permissive permettant son utilisation dans des projets commerciaux comme non commerciaux. Elle est notamment utilisée ou développée par des acteurs majeurs du secteur tels qu'Airbus Defence and Space, Thales Alenia Space, le CNES, l'ESA ou encore l'U.S. Naval Research Laboratory.

Orekit fournit principalement des fonctions de bas niveau, conçues pour être assemblées et adaptées en fonction de l'application considérée. La bibliothèque couvre ainsi un large éventail de domaines : la dynamique du vol spatial, la propagation et la détermination d'orbite, la gestion des attitudes, des dates et des systèmes de coordonnées, ou encore le traitement de mesures.


\section{Problématique générale}

En ce qui concerne la navigation par satellite, Orekit dispose de nombreuses fonctionnalités permettant de manipuler les données issues des signaux satellitaires. La bibliothèque ne fournit cependant pas directement de solution complète permettant de réaliser un positionnement par satellite. Plusieurs briques élémentaires sont déjà présentes, mais leur intégration au sein d'une chaîne complète reste à réaliser.

L'objectif du stage est donc de déterminer dans quelle mesure les fonctionnalités existantes d'Orekit permettent de construire une telle chaîne, puis d'identifier et de développer les éléments manquants afin d'atteindre progressivement des méthodes de positionnement de plus en plus précises.

\section{État de l'art}

Cette section présente un état de l'art des méthodes existantes de positionnement par satellite et établit un bilan des fonctions offertes par Orekit pour ce type d'application.

\subsection{GNSS (\textit{Global Navigation Satellite System})}

Le terme GNSS (\textit{Global Navigation Satellite System}) désigne l'ensemble des systèmes mondiaux de navigation par satellite permettant à un utilisateur équipé d'un récepteur de déterminer sa position, sa vitesse et une référence temporelle à partir des données issues des signaux émis par des satellites.

Les principaux systèmes actuellement opérationnels sont le GPS (\textit{Global Positioning System}) américain, GALILEO européen, GLONASS russe et BEIDOU chinois. Chacun repose sur une constellation de satellites placés en orbite autour de la Terre et diffusant en continu des signaux de navigation dans plusieurs bandes de fréquences. Ces signaux contiennent des messages de navigation permettant au récepteur d'exploiter les informations nécessaires au calcul de sa position. 

Les satellites GNSS émettent dans plusieurs bandes de fréquences situées dans
la bande L. Le GPS utilise par exemple les fréquences L1 à
$1575{,}42~\mathrm{MHz}$, L2 à $1227{,}60~\mathrm{MHz}$ et L5 à
$1176{,}45~\mathrm{MHz}$. Galileo émet notamment sur E1 à
$1575{,}42~\mathrm{MHz}$, E5a à $1176{,}45~\mathrm{MHz}$, E5b à
$1207{,}14~\mathrm{MHz}$ et E6 à $1278{,}75~\mathrm{MHz}$. Certaines
fréquences sont donc communes à plusieurs systèmes, comme GPS L1 et Galileo E1.

Chaque signal repose sur une porteuse radiofréquence modulée par une ou
plusieurs séquences pseudo-aléatoires, appelées codes PRN
(\textit{Pseudo-Random Noise}). Ces séquences sont connues du récepteur. Elles
permettent notamment d'identifier et de suivre les signaux reçus.

Pour le GPS et Galileo, plusieurs satellites d'une même constellation peuvent
émettre simultanément dans une même bande de fréquences. Leur séparation
repose notamment sur le CDMA (\textit{Code Division Multiple Access}) :
les signaux sont associés à des codes PRN différents, choisis pour présenter
de faibles corrélations croisées. Le récepteur génère localement les codes
attendus et effectue leur corrélation avec le signal reçu. Un pic de
corrélation permet de détecter le signal correspondant et d'estimer le
décalage temporel du code reçu par rapport à sa réplique locale.


\subsection{Principe de la mesure de distance}

Le positionnement GNSS repose sur l'estimation du temps de propagation des
signaux entre plusieurs satellites et le récepteur. Dans un cas idéal, si le
satellite $s$ émet un signal à l'instant $t_e^s$ et que celui-ci est reçu à
l'instant $t_r$, la distance parcourue est

\begin{align}
\rho_r^s = c\left(t_r-t_e^s\right),
\label{eq:geometric_time_range}
\end{align}

où $\rho_r^s$ est la distance géométrique entre le satellite et le récepteur et
$c$ la vitesse de la lumière dans le vide.

Cette relation conduit au principe de trilatération. Connaître la distance entre 
un satellite et le récepteur permet de déterminer que le récepteur se situe quelque 
part sur une sphère centrée sur ce satellite, dont le rayon est égal à la distance 
mesurée. Dans un cas idéal,
trois distances indépendantes permettent de déterminer les trois coordonnées
du récepteur, en écartant la solution qui ne se trouve pas au voisinage de la
surface terrestre.

En pratique, l'horloge du récepteur n'est pas parfaitement synchronisée avec
le temps du système GNSS. Son décalage constitue une quatrième inconnue,
commune aux mesures réalisées à une même époque. Il faut donc disposer des
mesures d'au moins quatre satellites pour déterminer les trois coordonnées du
récepteur et son biais d'horloge. Lorsque davantage de satellites sont
disponibles, le système devient surdéterminé et la redondance des observations
peut être exploitée lors de l'estimation.


\section{Observations GNSS}

Un récepteur GNSS fournit principalement deux types d'observations utilisables
pour le positionnement : les mesures de pseudo-distance issues du suivi des
codes et les mesures de phase de la porteuse. Par la suite, 
les termes codes ou pseudo-distance seront considérés comme équivalents.




\subsection{SPP (\textit{Single Point Positioning})}

Le SPP (\textit{Single Point Positioning}) est une méthode de positionnement
reposant principalement sur les mesures de pseudo-distance. Le
récepteur détermine sa position à partir des signaux reçus
et des informations de navigation diffusées par les satellites.

À chaque époque, les trois coordonnées du récepteur et son biais d'horloge sont
estimés à partir des observations disponibles. Les principales erreurs liées
aux horloges et à la propagation du signal sont corrigées à l'aide des
informations diffusées par les satellites et de modèles adaptés.

Le SPP ne nécessite ni station de référence ni produits externes précis. Sa
précision reste cependant limitée par la précision des données diffusées et
par les erreurs résiduelles des mesures et des modèles. Dans de bonnes
conditions de réception, la précision obtenue est généralement de l'ordre de
quelques mètres. Le SPP utilise une méthode de résolution des équations par moindres carrés pondérés (\textit{Weighted Least Squares}, WLS), qui sera détaillée plus tard.


\subsection{PPP (\textit{Precise Point Positioning})}

Le PPP (\textit{Precise Point Positioning}) est également une méthode de
positionnement. Il
cherche cependant à atteindre une précision supérieure au SPP en exploitant
des produits précis pour les satellites, les mesures de phase de la porteuse
et des modèles de correction plus complets.

L'utilisation de la phase améliore fortement la précision des observations,
mais impose d'estimer les ambiguïtés de phase. Le PPP s'appuie
également sur des données précises concernant les orbites et les horloges des
satellites afin que leur incertitude ne limite pas l'exploitation des mesures
de phase.

La solution PPP nécessite généralement une phase de convergence pendant
laquelle les différents paramètres sont estimés. Une fois cette convergence
atteinte, la précision peut atteindre quelques centimètres.

